% Generated draft: requires mathematical and source review.
\documentclass[11pt]{article}
\usepackage[margin=25mm]{geometry}
\usepackage{fontspec}
\setmainfont{Latin Modern Roman}
\usepackage{amsmath,amssymb,mathtools,unicode-math}
\setmathfont{Latin Modern Math}
\usepackage{enumitem,needspace,xurl,hyperref}
\setlength{\parindent}{0pt}
\setlength{\parskip}{5pt}
\setlength{\emergencystretch}{3em}
\newcommand{\N}{\mathbb{N}}
\newcommand{\Z}{\mathbb{Z}}
\newcommand{\Q}{\mathbb{Q}}
\newcommand{\R}{\mathbb{R}}
\newcommand{\C}{\mathbb{C}}
\newcommand{\F}{\mathbb{F}}
\newcommand{\A}{\mathbb{A}}
\newcommand{\PP}{\mathbb{P}}
\newcommand{\E}{\mathbb{E}}
\newcommand{\eps}{\varepsilon}
\newcommand{\dd}{\,\mathrm{d}}
\begin{document}
\section*{UnsolvedMath ID: 30000174}
\textbf{Source problem code:} catalogue entry
\subsection{Definitions and mathematical statement}
Let $S_g$ be the closed orientable surface of genus $g\ge0$. For an embedded graph $G$, edgewidth $\operatorname{ew}(G)$ is the minimum length of a noncontractible cycle, and facewidth $\operatorname{fw}(G)$ is the minimum number of intersections with $G$ of a noncontractible simple closed curve in general position. Use infinity if no such curve exists, as on the sphere. The dual $G^*$ requires the source's cellular-embedding convention. A cubic graph is three-edge-colorable if its edges can be colored with three colors so that edges incident at a vertex have different colors.

The source preserves three nested assertions:
\begin{enumerate}
\item For every $g\ge0$, $\chi_w(S_g)=4$.
\item More strongly, for every graph $G$ embedded in $S_g$ with $\operatorname{ew}(G^*)\ge3$, the printed assertion is $\phi_c(G)=4$. The source calls this equivalent to three-edge-colorability of every cubic embedded $G$ with $\operatorname{fw}(G)\ge3$ and reports $\operatorname{fw}(G)=\operatorname{ew}(G^*)$ for its cubic embeddings.
\item For every $g\ge0$ there exists $w$ such that every cubic graph embedded in $S_g$ with $\operatorname{ew}(G)\ge w$ is three-edge-colorable.
\end{enumerate}
The last assertion uses primal edgewidth, not dual edgewidth or facewidth. The original passage's informal phrase ``snarks on $S_g$ have bounded facewidth'' does not by itself justify replacing that displayed condition by another width.

Here $\phi_c$ conventionally denotes circular flow number: the infimum of $r\ge2$ for which an orientation supports a real flow, conserved at every vertex, with every directed edge value in $[1,r-1]$. Under this convention, three-edge-colorability of a bridgeless cubic graph corresponds to a nowhere-zero four-flow and $\phi_c(G)\le4$, not to equality with four. Smaller flow numbers are possible. The literal equality and asserted equivalence therefore require reconciliation; the sign is not silently changed. The excerpt also omits the admissible graph/connectivity hypotheses, which matter for flow existence.

The definition of the chromatic-width invariant $\chi_w$ appears in the missing source display (3). The original passage describes it through circular chromatic numbers $\chi_c$ of sufficiently wide embeddings and claims the equivalent orientable-surface target
\[
 \forall g\ \forall\varepsilon>0\ \exists W\ \forall G\hookrightarrow S_g,
 \quad\operatorname{ew}(G)\ge W\Longrightarrow\phi_c(G)<4+\varepsilon,
\]
with reduction to cubic graphs by lifting [7]. It cites [2] for $\chi_c(G)\le\chi_{\mathrm{loc}}(G)+\varepsilon$ at sufficiently large width and refers to missing display (4). Definitions and hypotheses for $\chi_{\mathrm{loc}}$ and those equivalences are not supplied.

The original historical discussion reports Fisk triangulations with exactly two odd-degree vertices, adjacent to each other, having $\chi(G)\ge5$ [3], and examples of arbitrarily large edgewidth for each positive genus. It gives only an unpublished $\chi_{\mathrm{loc}}<5$ result for a very special torus subclass, and a conditional near-four circular-coloring consequence of assertion 3; these are not universal proofs. It attributes the stronger conjecture to Gr\"unbaum [4].

The literature assessment checked on 21 August 2026 instead describes surface Pfaffian equivalences and cites a Pfaffian-orientation survey. Its fit to these circular-coloring/flow conjectures is unclear. That mismatch, missing definitions and the equality issue prevent treating its open label as an independently verified status conclusion.
\subsection{Short English statement}
The source asks three all-genus width and coloring conjectures, culminating in three-edge-colorability of sufficiently wide cubic embeddings. Its literal circular-flow equality, omitted graph conventions and mismatched Pfaffian literature assessment must be checked before the asserted equivalences can be accepted.
\subsection{Sources}
\begin{itemize}
\item \emph{Graph Theory} (2005), Oberwolfach Reports, DOI 10.4171/owr/2005/03, pp. 135--183, \url{https://doi.org/10.4171/owr/2005/03}. Original passage cites [2] for the local/circular comparison, [7] for lifting, Gr\"unbaum [4], and Steve Fisk [3]; full bibliography entries and displays (3), (4) are not supplied.
\item Literature assessment checked 21 August 2026: Norine, ``survey of Pfaffian orientations,'' \url{https://thomas.math.gatech.edu/PAP/pfafsurv.pdf}; original report, \url{https://doi.org/10.4171/owr/2005/03}. The assessment's Pfaffian focus is not reconciled with the corrected target.
\item Local UnsolvedMath registry, permanent ID $30000174$, exact source code \texttt{catalogue entry}, created 25 August 2026 and updated 16 September 2026. The clean statement is labelled corrected and verified with all three conjectures retained; the full original discussion supplies the historical qualifications recorded above.
\end{itemize}
\end{document}
